\chapter{Происхождение физического ядра из родительского действия}
\label{ch:v8-physical-correlation-kernel-parent-action-origin-gate}

\section{Может ли действие само определить процесс}

Предыдущий гейт показал: если полная матрица $\mathcal C_\tau$ известна
независимо, она восстанавливает относительные скорости. Теперь проверим
обратную стрелку: достаточно ли уже построенного родительского действия,
чтобы однозначно получить эту матрицу.

В квадратичном приближении действие задаёт положительный гессиан $H$ и
равновесную гауссовскую ковариацию
\begin{equation}
 \Sigma=H^{-1}.
 \label{eq:v8-parent-equilibrium-covariance}
\end{equation}
Однако градиентная динамика требует ещё положительной мобильности $M$:
\begin{equation}
 A=MH,
 \qquad
 \mathcal C_\tau=e^{-\tau A},
 \qquad
 A\Sigma+\Sigma A^{*}=2M.
 \label{eq:v8-parent-mobility-identity}
\end{equation}
Последнее тождество означает, что одна и та же мера с ковариацией $H^{-1}$
стационарна для любого $M>0$. Действие определяет рельеф, но не скорость
движения по нему.

\section{Контроль на фактическом полярном блоке}

В качестве наиболее благоприятного теста первые четыре семейных жёсткости
положены единичными, а в сектор $(QLYR,XLdR)$ вставлен уже выведенный
полярный блок
\begin{equation}
 H_{QX}=
 \begin{pmatrix}
 0.92492997&-0.44926187\\
 -0.44926187&0.58178422
 \end{pmatrix}.
 \label{eq:v8-parent-qx-hessian}
\end{equation}
Для этого единственного $H$ проверены пять мобильностей: изотропная,
семейно-анизотропная, ускоряющая сектор переноса, ускоряющая калибровочный сектор и
расщепляющая две межсекторные стрелки. Все пять точно удовлетворяют
\eqref{eq:v8-parent-mobility-identity}: максимальный остаток меньше
$1.84\cdot10^{-15}$.

Тем не менее их ядра при $\tau=0.4$ различны. Попарные расстояния Фробениуса
лежат между $0.4516$ и $1.0782$, а нормированные коротковременные скорости
варьируются от равномерного вектора $(1/6,\ldots,1/6)$ до, например,
\begin{equation}
 \frac1{15}(4,1,1,1,4,4)
 \quad\hbox{и}\quad
 \frac1{15}(1,4,4,4,1,1).
 \label{eq:v8-parent-distinct-mobilities}
\end{equation}
Следовательно, даже идеально известный квадратичный родитель не выбирает
шесть интенсивностей без дополнительной динамической метрики.

\section{Почему прежняя межсекторная ковариация не закрывает разрыв}

Для одного и того же $H_{QX}$ три ранее допустимых правила
$H^{-1}$, $H^{-1/2}$ и $e^{-H}$ дают отношения диагональных компонент
$QLYR:XLdR$, равные соответственно
\begin{equation}
 0.91681,qquad 0.95737,qquad 0.58000,
\end{equation}
и разные внедиагональные корреляции. Полярный гессиан устойчиво выбирает
главную ось, но не меру, не мобильность и не физический временной масштаб.

Аудит покрытия остальных семейств усиливает вывод. Связывающая инцидентность
задаёт форму, но её коэффициент остаётся свободным. Представления
$SU(3)$, $SU(2)$ и $U(1)$ известны, однако общий физический калибровочный след и
пространственно-временная кинетическая метрика в Томе VII не получены.
Ранняя TOE-программа, в свою очередь, постулирует гауссово ядро; мост с
тепловым ядром записан лишь как приближённая гипотеза.

\begin{theorem}[Разделение равновесия и динамики]
Родительский гессиан $H$ определяет гауссовскую равновесную ковариацию
$H^{-1}$, но не определяет единственное физическое корреляционное ядро.
Каждая положительная мобильность $M$ задаёт иной процесс $e^{-\tau MH}$ с
той же стационарной мерой. Поэтому существующее действие Тома VII не
выводит шесть скоростей Тома VIII.
\end{theorem}

\begin{nogocriterion}[Запрет отождествления гессиана с генератором]
Нельзя без отдельного принципа положить $M=I$, объявить
$\mathcal C_\tau=e^{-\tau H}$ и считать ядро физически выведенным.
Единичная мобильность является допустимым представителем, но столь же
допустимы положительные семейно-зависимые мобильности. Требуется вывести
закон флуктуации--диссипации, спектр средыальную плотность или общий
кинетический след.
\end{nogocriterion}

\begin{protocol}
\begin{enumerate}
 \item фактический полярный межсекторный гессиан: \textbf{использован};
 \item положительных мобильностей в контроле: \textbf{5};
 \item общая стационарная ковариация: \textbf{точно одна и та же};
 \item максимальный остаток стационарности: \textbf{$1.84\cdot10^{-15}$};
 \item попарно различных конечновременных ядер: \textbf{10 из 10};
 \item минимальное расстояние ядер: \textbf{$0.4516$};
 \item правило межсекторной ковариации: \textbf{не единственно};
 \item физическая калибровочная мобильность: \textbf{не получена};
 \item физическое $\mathcal C_\tau$ из существующего действия:
 \textbf{не получено};
 \item следующий гейт: \textbf{происхождение общей
 флуктуационно-диссипативной мобильности}.
\end{enumerate}
\end{protocol}

Машинный сертификат сохранён в
\path{s2t_v8_physical_correlation_kernel_parent_action_origin_gate_results.json}.

\begin{thebibliography}{9}
\bibitem{v8-kernel-parent-os}
K.~Osterwalder, R.~Schrader,
\emph{Axioms for Euclidean Green's Functions},
Commun. Math. Phys. 31 (1973) 83--112.
\bibitem{v8-kernel-parent-davies}
E.~B.~Davies,
\emph{Markovian Master Equations},
Commun. Math. Phys. 39 (1974) 91--110.
\bibitem{v8-kernel-parent-parisi-wu}
G.~Parisi, Y.-S.~Wu,
\emph{Perturbation Theory Without Gauge Fixing},
Sci. Sin. 24 (1981) 483--496.
\end{thebibliography}